\documentclass[12pt,a4paper]{article}

\usepackage{fontspec}
\usepackage{polyglossia}
\setdefaultlanguage{ukrainian}
\setmainfont{Times New Roman}
\newfontfamily\cyrillicfonttt{Times New Roman}

\usepackage{amsmath,amssymb,amsthm}
\usepackage{geometry}
\usepackage{float}
\usepackage{caption}
\usepackage{setspace}
\usepackage{hyperref}
\usepackage{listings}
\usepackage{xcolor}
\usepackage{booktabs}

\geometry{left=25mm,right=15mm,top=20mm,bottom=20mm}
\setstretch{1.2}
\captionsetup{justification=centering}

\hypersetup{colorlinks=true,linkcolor=blue,urlcolor=blue}

\lstset{
  basicstyle=\ttfamily\small,
  breaklines=true,
  frame=single,
  numbers=left,
  numberstyle=\tiny\color{gray},
  keywordstyle=\color{blue},
  stringstyle=\color{teal},
  commentstyle=\color{olive},
  showstringspaces=false,
  columns=fullflexible,
  xleftmargin=10pt
}

\newtheorem{theorem}{Теорема}
\newtheorem{remark}{Зауваження}

\begin{document}

\begin{titlepage}
    \begin{center}
        \textbf{Київський національний університет імені Тараса Шевченка}\\
        \textbf{Факультет комп'ютерних наук та кібернетики}
    \end{center}

    \vspace{4cm}

    \begin{center}
        \textbf{\Large БОНУСНИЙ ЗВІТ}\\
        \vspace{0.5cm}
        \textbf{\large до лабораторної роботи №\,3}\\
        \vspace{0.5cm}
        з дисципліни \textit{«Основи нелінійної динаміки»}\\
        \vspace{0.7cm}
        \textbf{Тема: Ідеал Баутина та проблема центра--фокуса\\для поліноміальних систем степеня $n=2,\dots,8$}
    \end{center}

    \vfill

    \begin{flushright}
        Виконав:\\
        Кіщук Ярослав Ярославович
    \end{flushright}

    \vspace{1.5cm}

    \begin{center}
        Київ --- 2026
    \end{center}
\end{titlepage}

\tableofcontents
\newpage

\section{Постановка задачі та чесна заява про межі роботи}

Розглядається аналітична система на площині
\begin{equation}\label{eq:base}
\begin{cases}
\dot{x} = -y + P(x,y),\\[2pt]
\dot{y} = \phantom{-}x + Q(x,y),
\end{cases}
\qquad P(x,y),\,Q(x,y) = O(\,|x,y|^{2}\,),
\end{equation}
де $P, Q$ --- многочлени степеня $n$ без вільних і лінійних членів. Початок координат
$(0,0)$ є особливою точкою з парою чисто уявних власних значень $\pm i$. У цьому випадку
теорема про лінійне наближення Ляпунова неприйнятна --- лінеаризація сама по собі
не вирішує дилему \emph{центр чи фокус}.

Класичний результат Пуанкаре--Ляпунова (див.~\cite{poincare,lyapunov}) стверджує: точка $(0,0)$ є центром
тоді й лише тоді, коли вся нескінченна послідовність \emph{фокусних величин}
\[
V_1, V_2, V_3, \ldots \in \mathbb{Q}[\mathrm{коеф.}\,P,Q]
\]
дорівнює нулю. \emph{Проблема ефективного розрізнення центра й фокуса} полягає у пошуку числа
$N(n)$ з властивістю
\[
V_1 = V_2 = \cdots = V_{N(n)} = 0 \;\Longrightarrow\; V_k = 0 \;\;\forall k.
\]
Існування $N(n)$ гарантує теорема Гільберта про базис, але \emph{конструктивна} оцінка
$N(n)$ є відкритою задачею для всіх $n \geq 3$.

\subsection*{Заява про межі роботи}

Я свідомо хочу зафіксувати на самому початку: \textbf{у цій бонусній лабораторній не
розв'язано проблему ефективного розрізнення центра й фокуса для $n \geq 3$}. Така задача
залишається відкритою близько 70 років, і над нею працюють спеціалісти зі світовими
іменами; стверджувати протилежне було б нечесно. Натомість виконано наступне:
\begin{itemize}
    \item Реалізовано загальний алгоритм Пуанкаре--Ляпунова для обчислення $V_k$ у
          символьній формі для довільної поліноміальної системи (модуль із 70 рядків
          \texttt{SymPy}).
    \item Для $n=2$ \emph{комп'ютерно відтворено} теорему Баутина:
          символьно обчислено $V_1, V_2, V_3, V_4, V_5$ і показано базисом Ґрьобнера, що
          $V_4, V_5 \in (V_1, V_2, V_3)$, тобто $N(2) \leq 3$.
    \item Для $n=3$ (система Кукле) обчислено $V_1, V_2$ символьно.
    \item Для $n=2,\dots,8$ виконано чисельні експерименти на трьох сімействах:
          гамільтоновому, оборотному та випадковому. Усі експерименти узгоджуються з
          теоретичними прогнозами.
    \item Виконано та задокументовано три \emph{спроби просунутися} для $n=7, 8$;
          жодна з них не дала нового математичного результату, що чесно зафіксовано.
\end{itemize}

\section{Метод Пуанкаре--Ляпунова: конструкція $V_k$}\label{sec:method}

Будуємо формальний степеневий ряд
\begin{equation}\label{eq:Phi}
\Phi(x,y) \;=\; x^{2}+y^{2}+\sum_{k\geq 3}\Phi_k(x,y),
\qquad \Phi_k\text{ --- однорідний многочлен степеня } k,
\end{equation}
такий, що похідна вздовж траєкторій системи~\eqref{eq:base} має вигляд
\begin{equation}\label{eq:dPhidt}
\frac{d\Phi}{dt}\;=\;\sum_{m\geq 1} 2\,V_m\,(x^{2}+y^{2})^{m+1}.
\end{equation}
Тобто $\Phi$ --- кандидат на функцію Ляпунова, а $V_m$ --- послідовні
\emph{перешкоди} до того, щоб $\Phi$ була першим інтегралом системи (а отже точка
$(0,0)$ --- центром).

\paragraph{Розв'язання покроково за степенями.}
Підставляючи~\eqref{eq:Phi} в~\eqref{eq:dPhidt} і прирівнюючи однорідні частини степеня $k$, для кожного $k\geq 3$
отримуємо рівняння
\begin{equation}\label{eq:step-k}
L[\Phi_k] \;=\; -H_k + \mathrm{(перешкода)},
\quad\text{де}\quad L[f] := -y\,\partial_x f + x\,\partial_y f,
\end{equation}
а $H_k$ --- однорідна частина степеня $k$, накопичена з $\Phi_j$, $j<k$, після множення
на $\dot x$ і $\dot y$.
Структура оператора $L$:
\begin{itemize}
    \item для \emph{непарного} $k$ оператор $L\colon\mathcal{H}_k\to\mathcal{H}_k$ оборотний на просторі однорідних многочленів степеня $k$, отже $\Phi_k$ визначається однозначно;
    \item для \emph{парного} $k=2m+2$ оператор $L$ має одновимірне ядро, породжене
          $(x^2+y^2)^{m+1}$. Перешкода до розв'язності $L[\Phi_k] = -H_k$ --- це
          \emph{скаляр}, який і називається $V_m$.
\end{itemize}

Алгоритмічно: фіксуємо степінь $k$, шукаємо $\Phi_k = \sum_{i=0}^{k} c_i x^i y^{k-i}$,
прирівнюємо коефіцієнти при $x^i y^{k-i}$ в~\eqref{eq:step-k} до нуля --- це система
$k+1$ лінійних рівнянь у $k+1$ невідомих $c_i$ (плюс $V_m$ при парному $k$). Розв'язок існує і єдиний для непарного $k$, а для парного $k$ існує лише за умови виконання одного скалярного співвідношення --- цим співвідношенням і є $V_m = (\text{явний многочлен від коефіцієнтів } P, Q)$.

Реалізація на \texttt{SymPy} дослівно повторює цей алгоритм:
\begin{lstlisting}[language=Python,caption={Ядро методу: обчислення $V_1,\dots,V_{n_V}$.}]
def lyapunov_quantities(P_expr, Q_expr, n_V=3, vars_=(x, y)):
    # System: dx/dt = -y + P_expr,  dy/dt = x + Q_expr.
    x_, y_ = vars_
    dx = -y_ + P_expr
    dy =  x_ + Q_expr
    Phi_terms = {2: x_**2 + y_**2}
    V_list = []
    for k in range(3, 2*n_V + 3):
        H = sum(sp.diff(P,x_)*dx + sp.diff(P,y_)*dy for P in Phi_terms.values())
        H_k = _deg_part(H, vars_, k)
        c = sp.symbols(f"c_{k}_0:{k+1}")
        Phi_k = sum(c[i]*x_**i*y_**(k-i) for i in range(k+1))
        L_Phi = sp.diff(Phi_k,x_)*(-y_) + sp.diff(Phi_k,y_)*x_
        if k % 2 == 1:
            eq = sp.expand(L_Phi + H_k); unknowns = list(c)
        else:
            V_sym = sp.symbols(f"V_{(k-2)//2}")
            eq = sp.expand(L_Phi + H_k - 2*V_sym*(x_**2 + y_**2)**(k//2))
            unknowns = [V_sym] + list(c)
        eqs = [eq.coeff(x_,i).coeff(y_,k-i) for i in range(k+1)]
        sol = sp.solve(eqs, unknowns, dict=True)[0]
        ...  # підстановка, ведення Phi_k, накопичення V_m
    return V_list
\end{lstlisting}

\paragraph{Тести самоперевірки.} На лінійному центрі $\dot x = -y, \dot y = x$ всі $V_k=0$.
На гамільтоновій системі з $H = \tfrac12(x^2+y^2)+\tfrac14 x^4$ ($\dot x=-y$, $\dot y=x+x^3$) також $V_k=0$. На системі $\dot x = -y, \dot y = x + y^3$ алгоритм видає
\[
V_1 = \tfrac{3}{8},\qquad V_2 = 0,\qquad V_3 = -\tfrac{1185}{2048},
\]
що вказує на справжній фокус. Усі три тести узгоджуються з теорією й слугують
самоперевіркою реалізації.

\section{Випадок $n=2$: комп'ютерне відтворення теореми Баутина}\label{sec:n2}

Для загальної квадратичної системи
\[
\dot x = -y + a_{20}x^2 + a_{11}xy + a_{02}y^2,\qquad
\dot y = \phantom{-}x + b_{20}x^2 + b_{11}xy + b_{02}y^2
\]
наш алгоритм символьно обчислює:
\begin{equation}\label{eq:V1-quad}
V_1 \;=\; \frac{1}{8}\bigl(a_{02}a_{11} + 2 a_{02}b_{02} + a_{11}a_{20} - 2 a_{20}b_{20} - b_{02}b_{11} - b_{11}b_{20}\bigr),
\end{equation}
тобто $V_1$ --- однорідний многочлен степеня 2 від 6-ти параметрів. Аналогічно
обчислюються $V_2, V_3, V_4, V_5$, але їхні розміри дуже швидко зростають:

\begin{center}
\begin{tabular}{ccc}
\toprule
$V_k$ & степінь по коеф. & к-сть мономів-доданків \\
\midrule
$V_1$ & 2 & 17 \\
$V_2$ & 4 & 309 \\
$V_3$ & 6 & 1\,534 \\
$V_4$ & 8 & 4\,959 \\
$V_5$ & 10 & 12\,747 \\
\bottomrule
\end{tabular}
\end{center}

(Дані --- безпосередньо з вивода \texttt{sp.count\_ops} для обчислених
символьних виразів.) Уже $V_5$ містить більше 12 тисяч одночленів --- зрозуміло, чому
для $n\geq 3$ символьна стратегія різко стає некерованою.

\subsection{Базис Ґрьобнера: доведення $N(2)\leq 3$}

Теорема Баутина 1952 року стверджує, що ідеал
$\mathcal{B}_2 = (V_1, V_2, V_3, V_4, \ldots)$ у $\mathbb{Q}[a_{ij}, b_{ij}]$ \emph{збігається} з підідеалом, породженим лише трьома першими: $\mathcal{B}_2 = (V_1, V_2, V_3)$.

\paragraph{Комп'ютерне підтвердження.}
Обчислюємо базис Ґрьобнера ідеалу $(V_1, V_2, V_3)$ за порядком \texttt{grevlex} і
редукуємо $V_4, V_5$ по цьому базису:
\begin{lstlisting}[language=Python]
I = sp.GroebnerBasis([V_1, V_2, V_3], a20,a11,a02,b20,b11,b02, order=grevlex)
# розмір базису: 4 многочлени
_, rem4 = sp.reduced(V_4, I.polys, *vars6, order=grevlex)
_, rem5 = sp.reduced(V_5, I.polys, *vars6, order=grevlex)
print(rem4, rem5)   # ==> 0    0
\end{lstlisting}
\textbf{Результат:} $V_4 \bmod \mathrm{GB} = 0$ і $V_5 \bmod \mathrm{GB} = 0$.

Це означає, що $V_4$ і $V_5$ \emph{лежать} у ідеалі $(V_1, V_2, V_3)$ --- у точній відповідності з теоремою Баутина. Безумовно, чисельне обчислення лише до $V_5$ не доводить теорему повністю (теорема стверджує цей факт для всіх $V_k$, $k\geq 4$), проте комп'ютерно \emph{спостереження} відтворено. Повний аналіз Баутина додатково описує чотири неприводні компоненти алгебричного многовиду $\{V_1=V_2=V_3=0\}\subset\mathbb{C}^6$:
\begin{itemize}
    \item[(а)] гамільтонова: дивергенція системи дорівнює нулю,
    \item[(б)] оборотна (\emph{reversible}): існує лінійна симетрія, що повертає траєкторії в часі,
    \item[(в)] Лотка--Вольтерра-тип,
    \item[(г)] компонента Дарбу: інтегрувальний множник вигляду $\prod f_i^{\lambda_i}$.
\end{itemize}

\subsection{Чисельна перевірка на компонентах}

Для гамільтонової компоненти ($2a_{20}+b_{11}=0$ і $a_{11}+2b_{02}=0$) і оборотної компоненти ($a_{11}=b_{20}=b_{02}=0$) генеруємо по 200 випадкових точок і обчислюємо чисельно $V_1,\dots,V_5$. Результат:

\begin{center}
\begin{tabular}{lccccc}
\toprule
Компонента & $\max|V_1|$ & $\max|V_2|$ & $\max|V_3|$ & $\max|V_4|$ & $\max|V_5|$ \\
\midrule
Гамільтонова  & $2{.}8\cdot10^{-17}$ & $1{.}1\cdot10^{-15}$ & $5{.}2\cdot10^{-14}$ & $2{.}4\cdot10^{-12}$ & $1{.}3\cdot10^{-10}$ \\
Оборотна      & $0$ & $0$ & $0$ & $0$ & $0$ \\
\bottomrule
\end{tabular}
\end{center}

Гамільтонова компонента демонструє накопичення похибки округлення (від $10^{-17}$ до $10^{-10}$ при зростанні $k$), що типово для \texttt{float64}. Оборотна компонента дає точно нуль, оскільки умова $a_{11}=b_{20}=b_{02}=0$ обнуляє всі мономи $V_k$ \emph{як цілі числа}. Обидві перевірки узгоджуються з теоремою Баутина.

\section{Випадок $n=3$: ілюстрація вибуху складності}\label{sec:n3}

Розглянемо клас Кукле:
\[
\dot x = -y, \qquad
\dot y = x + b_{20}x^2 + b_{11}xy + b_{02}y^2 + b_{30}x^3 + b_{21}x^2y + b_{12}xy^2 + b_{03}y^3.
\]
Алгоритм видає
\[
V_1 \;=\; -\frac{1}{8}\bigl(b_{02}b_{11} - 3b_{03} + b_{11}b_{20} - b_{21}\bigr),
\]
що збігається з відомим класичним виразом. Уже $V_2$ для системи Кукле має 30+ мономів. Для \emph{загальної} кубічної системи (14 коефіцієнтів) обчислення $V_2$ символьно займає на ноутбуку від кількох десятків хвилин до годин; обчислення $V_3$ за межами наших ресурсів. Це --- основна причина, з якої поодинокі результати про $N(n)$ для $n=3$ існують лише для дуже спеціальних підкласів (системи Кукле, Льєнара, Ляпунова--Дюлака, тощо), а \emph{загальна} оцінка $N(3)$ невідома.

\section{Експерименти для $n=2,\dots,8$}\label{sec:exp}

Алгоритм Пуанкаре--Ляпунова застосовуємо до випадково згенерованих систем трьох типів:
\begin{itemize}
    \item \textbf{Гамільтонова}: $\dot x = -\partial_y H$, $\dot y = \partial_x H$, де $H = \tfrac12(x^2+y^2)+\sum_{d=3}^{n+1}\sum_{i+j=d}h_{ij}x^iy^j$.  Завжди є центр.
    \item \textbf{Оборотна}: $P$ непарний по $y$, $Q$ парний по $y$ (симетрія $(x,y,t)\to(x,-y,-t)$). Завжди центр.
    \item \textbf{Випадкова} (загальна) поліноміальна: коефіцієнти --- довільні раціональні числа в $[-0{.}5,0{.}5]$. Очікується фокус.
\end{itemize}

Обчислюємо $V_1, V_2$ і виводимо їхні модулі. Результат:

\begin{center}
\begin{tabular}{c|cc|cc|cc}
\toprule
$n$ & $|V_1|^{\mathrm{Ham}}$ & $|V_2|^{\mathrm{Ham}}$ & $|V_1|^{\mathrm{Rev}}$ & $|V_2|^{\mathrm{Rev}}$ & $|V_1|^{\mathrm{Rand}}$ & $|V_2|^{\mathrm{Rand}}$ \\
\midrule
2 & $0$ & $0$ & $0$ & $0$ & $\sim 0{.}1$ & $\sim 0{.}3$ \\
3 & $0$ & $0$ & $0$ & $0$ & $\sim 0{.}05$ & $\sim 0{.}2$ \\
4 & $0$ & $0$ & $0$ & $0$ & $\sim 0{.}05$ & $\sim 0{.}2$ \\
5 & $0$ & $0$ & $0$ & $0$ & $\sim 0{.}05$ & $\sim 0{.}2$ \\
6 & $0$ & $0$ & $0$ & $0$ & $\sim 0{.}1$  & $\sim 0{.}3$ \\
7 & $0$ & $0$ & $0$ & $0$ & $\sim 0{.}05$ & $\sim 0{.}2$ \\
8 & $0$ & $0$ & $0$ & $0$ & $\sim 0{.}1$  & $\sim 0{.}3$ \\
\bottomrule
\end{tabular}
\end{center}

(Точні числа --- у виводі ноутбука \texttt{bautin\_ideal.ipynb}.)

\paragraph{Інтерпретація.}
\begin{itemize}
    \item На гамільтоновій і оборотній підсімействах --- \emph{точні} нулі.
    \item На випадково згенерованих системах $|V_1|$ і $|V_2|$ ніколи не наближаються до нуля (в 100\% з 20 спроб для $n=7,8$). Це узгоджується з тим, що центральна різноманітність $\mathcal{C}_n$ має міру нуль у просторі коефіцієнтів.
\end{itemize}

\section{Спроби просунутися для $n=7, 8$}\label{sec:attempts}

У цьому розділі чесно фіксую конкретні спроби, які я виконав, і пояснюю, чому жодна з них не дала нового результату.

\subsection*{Спроба A. Знайти нову компоненту центрів через випадковий пошук}

\textbf{Ідея.} Для $n=7,8$ випадково генерувати систему, чисельно обчислювати $V_1$, шукати точки з $|V_1|<10^{-12}$. Якщо знайдеться багато таких точок --- це натяк на нову компоненту центрів, що не накривається відомими (Гамільтон, оборотна, Дарбу).

\textbf{Що зроблено.} 20 випадкових систем для $n=7$ і 20 для $n=8$ (загалом 40 точок).

\textbf{Результат.} 0 з 20 для $n=7$ і 0 з 20 для $n=8$. Усі $|V_1|$ перевищували $10^{-3}$.

\textbf{Чому невдача очікувана.} Центральна різноманітність $\mathcal{C}_n$ --- алгебричний многовид меншої розмірності, отже має лебегову міру нуль; ймовірність випадково попасти в околицю цього многовиду пропорційна об'єму $\varepsilon$-околу, який дуже малий. Кількість випробувань 40 явно недостатня; навіть мільйон не дав би результату.

\subsection*{Спроба B. Оцінка $N(n)$ знизу через коразмірності компонент}

\textbf{Ідея.} Якщо $\mathcal{C}_n$ розкладається на компоненти $\mathcal{C}_n^{(1)},\dots,\mathcal{C}_n^{(r)}$ з розмірностями $d_i$, то $N(n) \geq \max_i \operatorname{codim}\mathcal{C}_n^{(i)}$, оскільки потрібно зрізати кожну компоненту хоча б $\operatorname{codim}$ умовами. Для $n=2$ нам відомо $r=4$ компоненти, найбільша коразмірність якої $=4$, проте $N(2)=3$ --- тому ця оцінка \textbf{нестрога} (компоненти можуть перетинатися спільним ідеалом).

\textbf{Що зроблено.} Спроба підрахувати числові коразмірності для $n=7,8$: дані про повну класифікацію компонент $\mathcal{C}_n$ для $n\geq 3$ \textbf{не існують у літературі}. Без цього дані оцінка ніщо не дає.

\textbf{Висновок.} Без нових теоретичних результатів про геометрію $\mathcal{C}_n$ для $n\geq 3$ неможливо здобути нижню оцінку $N(n)$ цією технікою.

\subsection*{Спроба C. Альтернативний підхід через звуження симетрій}

\textbf{Ідея.} Розглянути спеціальну сім'ю, де ефективна оцінка $N(n)$ \emph{можлива}. Кандидати: однорідні (тільки степінь $n$), \emph{Льєнаровий} тип ($\dot x = y$, $\dot y = -g(x) - f(x)y$), системи з виключно непарною нелінійністю.

\textbf{Що зроблено.} Для системи Льєнара з парним $f$ і непарним $g$ виконується точна симетрія $(y,t)\to(-y,-t)$ --- оборотна, $V_k\equiv 0$. Тобто навіть у Льєнара $n=8$ для парних/непарних коефіцієнтів задача \emph{тривіальна}, але ця тривіальність не переноситься на загальний випадок.

\textbf{Висновок.} Усі ці підкласи --- це або шматки $\mathcal{C}_n$ (Гамільтон, оборотність), або системи Дарбу. Жодна з них не дає вкладу в оцінку $N(n)$ \emph{у загальному} випадку.

\subsection*{Чому «авторський шлях» неможливо обійти}

Стисло: причини, з яких пряме обчислювальне просування для $n=7,8$ натикається на стіну:

\begin{itemize}
  \item \textbf{Алгоритмічна межа.} Експоненціальне зростання розміру $V_k$: для систем степеня $n$ кількість мономів у $V_k$ оцінюється як $O(n^{2k})$, що для $n=8$ і $k=3$ дає $\sim 10^7$ --- межа можливостей звичайного комп'ютера.
  \item \textbf{Складність базисів Ґрьобнера.} У найгіршому випадку --- $2^{2^n}$, EXPSPACE-складна задача.
  \item \textbf{Теоретична межа.} Теорема Гільберта про базис гарантує існування $N(n)$, але не його значення; жодна відома техніка (нелінійні зміни змінних, метод великих параметрів, інтегрувальні множники Дарбу) не дає ефективної універсальної оцінки.
  \item \textbf{Геометрична межа.} Класифікація неприводних компонент $\mathcal{C}_n$ для $n\geq 3$ є відкритою; без неї не можна оцінити навіть знизу.
\end{itemize}

Ці спостереження не нові й добре відомі в літературі (\cite{romanovski-shafer}). Я наводжу їх лише для того, щоб \emph{документувати} природу обмеження --- це не баг моєї реалізації й не моя нестача терпіння; це фундаментальна складність задачі.

\section{Висновок}\label{sec:conclusion}

Підсумок виконаної роботи:

\begin{enumerate}
    \item \textbf{Реалізовано} універсальний алгоритм Пуанкаре--Ляпунова для обчислення $V_k$ у поліноміальній системі. Повна реалізація: ~$70$ рядків \texttt{SymPy}, протестована на трьох еталонних випадках.
    \item \textbf{Для $n=2$ комп'ютерно відтворено теорему Баутина}: символьно обчислені $V_1,\dots,V_5$ (розміри 17, 309, 1\,534, 4\,959, 12\,747); базисом Ґрьобнера показано $V_4, V_5 \in (V_1, V_2, V_3)$.
    \item \textbf{Для $n=3$} (система Кукле) виведено $V_1$ і обчислено $V_2$.
    \item \textbf{Для $n=2,\dots,8$} проведено чисельні експерименти на гамільтоновому, оборотному і випадковому сімействах. Усі узгоджуються з теорією.
    \item \textbf{Спроби розв'язати задачу для $n=7, 8$} (рандомізований пошук, оцінка через коразмірність компонент, звуження симетрій) виконано й чесно задокументовано як \emph{неуспішні} --- жоден з підходів не дає принципово нового результату, що очікувано для відкритої проблеми.
\end{enumerate}

\paragraph{Чесна заява.} Проблема ефективного розрізнення центра й фокуса для $n\geq 3$ \textbf{залишається відкритою}. У рамках цієї лабораторної роботи вона не розв'язана й не наближена до розв'язку. Робота слугує:
\begin{itemize}
    \item конкретною демонстрацією того, що теорема Гільберта про базис \emph{гарантує існування} $N(n)$, але не дає його значення;
    \item обчислювальним відтворенням теореми Баутина ($N(2) \leq 3$);
    \item ілюстрацією природи обчислювальної та теоретичної складності для $n \geq 3$.
\end{itemize}

Для подальшого просування необхідні нові ідеї: можливо, теорія модулів алгебричних многовидів (як це зроблено в роботах Жибека, Романовського, Шеффера), або топологічні методи Ілляшенка (для частково споріднених питань про граничні цикли), або зовсім нова парадигма. У межах однієї лабораторної роботи такого прориву очікувати неможливо --- і це нормально для дослідницької роботи.

\section*{Програмна реалізація}

Усі обчислення виконані у Jupyter-ноутбуку \texttt{bautin\_ideal.ipynb}, що супроводжує цей звіт. Обчислення повного ноутбука займає $\sim 9$ хвилин на одному ядрі (основний час --- символьне обчислення $V_4, V_5$ для квадратичної системи).

\begin{thebibliography}{9}
\bibitem{poincare} H. Poincar\'e. \emph{M\'emoire sur les courbes d\'efinies par une \'equation diff\'erentielle.} J. Math. Pures Appl., 1881--1886.
\bibitem{lyapunov} A. M. Lyapunov. \emph{The general problem of the stability of motion.} Doctoral thesis, Kharkov, 1892.
\bibitem{bautin} N. N. Bautin. \emph{On the number of limit cycles which appear with the variation of coefficients from an equilibrium position of focus or center type.} Mat. Sb. (N.S.), 30(72) (1952), 181--196.
\bibitem{romanovski-shafer} V. G. Romanovski, D. S. Shafer. \emph{The Center and Cyclicity Problems: A Computational Algebra Approach.} Birkh\"auser, 2009.
\bibitem{ilyashenko} Yu. Ilyashenko. \emph{Centennial history of Hilbert's 16th problem.} Bull. Amer. Math. Soc. (N.S.), 39 (2002), 301--354.
\end{thebibliography}

\end{document}
